% Wave-18 E1/20: the Atiyah flop as an automorphism of the positive half
% Yplus(\widetilde Y) of the resolved conifold.
% Voice: Bridgeland + Szendroi, channelled through Russian elite + math-physics elite.
% Primary sources: Bridgeland 2002 (arXiv:math/0502050); Szendroi 2008
% (arXiv:0705.3419); Kontsevich-Soibelman 2008 (arXiv:0811.2435) sec.8;
% Nagao-Nakajima 2011 (arXiv:0809.2992); Morrison-Mozgovoy-Nagao-Szendroi
% 2012 (arXiv:1107.5017); Van den Bergh 2004 (arXiv:math/0207170).
% Companion to wave17_c1_conifold_gluing_explicit.tex.

\providecommand{\C}{\mathbb{C}}
\providecommand{\Z}{\mathbb{Z}}
\providecommand{\N}{\mathbb{N}}
\providecommand{\bP}{\mathbb{P}}
\providecommand{\cO}{\mathcal{O}}
\providecommand{\cY}{\mathcal{Y}}
\providecommand{\cF}{\mathcal{F}}
\providecommand{\cE}{\mathcal{E}}
\providecommand{\Yplus}{Y^{+}}
\providecommand{\CoHA}{\mathrm{CoHA}}
\providecommand{\DT}{\mathrm{DT}}
\providecommand{\Db}{D^{b}}
\providecommand{\Coh}{\mathrm{Coh}}
\providecommand{\tr}{\mathrm{tr}}
\providecommand{\Jac}{\mathrm{Jac}}
\providecommand{\Aut}{\mathrm{Aut}}
\providecommand{\kch}{\kappa_{\mathrm{ch}}}
\providecommand{\kcat}{\kappa_{\mathrm{cat}}}
\providecommand{\kBKM}{\kappa_{\mathrm{BKM}}}
\providecommand{\slhat}{\widehat{\mathfrak{sl}}_{2}}
\providecommand{\nplus}{\mathfrak{n}_{+}^{\slhat}}
\providecommand{\ghat}{\widehat{\mathfrak{gl}}_{1}}
\providecommand{\Wflop}{\phi_{\mathrm{flop}}}
\providecommand{\Wflopp}{\phi_{\mathrm{flop}}^{+}}
\providecommand{\FMflop}{\Phi_{\mathrm{flop}}}
\providecommand{\ClaimStatusTheorem}{\textsc{[theorem]}}
\providecommand{\ClaimStatusDefinition}{\textsc{[definition]}}
\providecommand{\ClaimStatusDerivation}{\textsc{[first-principles]}}
\providecommand{\ClaimStatusConjectured}{\textsc{[conjectural]}}

\section*{Wave-18 E1/20: the Atiyah flop as an automorphism of
         $\Yplus(\widetilde Y)$}
\label{wn:sec:flop-mutation-Yplus}

The resolved conifold admits a second small resolution
$\widetilde{Y}^{\vee}$, obtained by contracting the exceptional
$\bP^{1}$ and blowing up along the transverse direction.  Bridgeland's
derived equivalence $\Db(\widetilde Y)\cong\Db(\widetilde Y^{\vee})$
lifts, through the cohomological-Hall construction, to an automorphism
$\Wflopp$ of the global positive half $\Yplus(\widetilde Y) =
U(\nplus)\otimes\Lambda$ of Wave-17 C1.  Five cycles show: the
automorphism is the affine Weyl reflection $s_{1}$ on
$A_{1}^{(1)}$, composed with a parity flip on the Cartan
$\Lambda$; its character-level shadow is the Kontsevich--Soibelman
wall-crossing of the conifold chamber; together with the overlap
cocycle $\phi_{UV}$ of Wave-17 C1 it generates the Seiberg-duality
groupoid of $(Q_{\mathrm{KW}},W_{\mathrm{KW}})$.

\subsection*{Cycle 1 --- the flop as Fourier--Mukai transform}

\ClaimStatusDefinition\ Let $\widetilde{Y}\times_{Y}\widetilde{Y}^{\vee}$
denote the fibre product over the singular conifold
$Y=\{xy=zw\}\subset\C^{4}$, and let $\cE = \cO_{\widetilde{Y}\times_{Y}
\widetilde{Y}^{\vee}}$ be its structure sheaf.  Bridgeland
(\texttt{math/0502050}, Thm.~1.1) shows that $\cE$ is a perfect complex
whose associated Fourier--Mukai kernel
\[
  \FMflop \;=\; \bigl(p_{2,*}\circ(\,\cdot\,\otimes^{\mathbb L}\cE)\circ
                   p_{1}^{*}\bigr)\colon \Db(\widetilde{Y})
                   \;\longrightarrow\; \Db(\widetilde{Y}^{\vee})
\]
is an equivalence of triangulated categories.  Away from the
exceptional $\bP^{1}$, $\FMflop$ restricts to the identity; on
the exceptional fibre it acts as the spherical twist around
$\cO_{\bP^{1}}(-1)$ (Seidel--Thomas-type twist).  The two resolutions
share the same toric fan at the level of rays but differ in
triangulation (Wave-17 C1, Cycle 1); the flop is the change of
triangulation.

\subsection*{Cycle 2 --- induced CoHA automorphism}

\ClaimStatusTheorem\ The Fourier--Mukai equivalence $\FMflop$ preserves
compact support, preserves the $(-1,-1)$-shifted symplectic structure
(the Calabi--Yau $3$-form, up to the orientation flip of Wave-17 C1
Cycle 2), and commutes with Ext pushforward along moduli of compactly
supported sheaves.  By the functoriality of the cohomological-Hall
construction for CY$_{3}$ categories (Kontsevich--Soibelman
\texttt{0811.2435} sec.~5; Davison--Meinhardt \texttt{1601.02479}
Thm.~A), $\FMflop$ induces an isomorphism of associative algebras
\[
  \Wflopp\colon \Yplus(\widetilde{Y})
  \;\xrightarrow{\,\sim\,}\; \Yplus(\widetilde{Y}^{\vee}).
\]
The two global positive halves are \emph{a priori} distinct algebras ---
one built from $(Q_{\mathrm{KW}},W_{\mathrm{KW}})$ with vertex set
$\{0,1\}$, the other with $\{0',1'\}$ and the opposite orientation of
each arrow.  Both are, by Wave-17 C1, Cycle 5, canonically isomorphic
to $U(\nplus)\otimes\Lambda$; the flop-induced map identifies the
two copies and lands as an endomorphism of the ambient algebra.

\subsection*{Cycle 3 --- explicit form: $s_{1}$-Weyl reflection and
             parity flip}

\ClaimStatusDerivation\ Under the identification
$\Yplus(\widetilde{Y})=U(\nplus)\otimes\Lambda$, the Chevalley
generators indexed by the $A_{1}^{(1)}$ simple roots $\alpha_{0}=\delta-\alpha,\alpha_{1}=\alpha$
are
\[
  e_{0} \;\longleftrightarrow\; \tr(a_{1}b_{2}-a_{2}b_{1}),\qquad
  e_{1} \;\longleftrightarrow\; \tr(b_{1}a_{2}-b_{2}a_{1}),
\]
i.e.\ the two real-root level-$(1,0)$ and $(0,1)$ generators of the MMNS
resolution (\texttt{1107.5017}, Prop.~2.4).  On the quiver side, Seiberg
duality at vertex $1$ exchanges $a_{i}\leftrightarrow b_{i}^{\vee}$,
reverses arrow orientation, and adjusts $W_{\mathrm{KW}}$ by the
dualising term; on the root system, this is the simple reflection
\[
  s_{1}\colon \alpha_{1}\mapsto -\alpha_{1},\qquad
  \alpha_{0}\mapsto \alpha_{0}+\alpha_{1},
\]
extended multiplicatively to the Kac--Moody Weyl group $W(A_{1}^{(1)})$.
Accordingly
\[
  \Wflopp(e_{0}) = e_{0} + [e_{1},e_{0}],\qquad \Wflopp(e_{1}) = -e_{1}
\]
(Serre relation: $\mathrm{ad}_{e_{1}}^{3}(e_{0})=0$ in $\slhat$, so the
formula closes at one bracket).  On the Cartan $\Lambda = \Lambda^{*}
(h_{\delta})$ (the Heisenberg generated by the imaginary-root elements),
the flop acts by \emph{parity flip} $h_{n\delta}\mapsto (-1)^{n}
h_{n\delta}$.  This is forced by the orientation-reversal of
$\Omega_{\widetilde{Y}}\mapsto -\Omega_{\widetilde{Y}^{\vee}}$ on the
overlap (Wave-17 C1, Cycle 2): each imaginary-root generator, being
built from $\bP^{1}$-fundamental classes, picks up the pullback sign of
the volume form.

The spectrum of invariants is preserved: $\kch(\Yplus(\widetilde Y)) =
\kch(\Yplus(\widetilde Y^{\vee})) = 1$, and $\kcat(\widetilde Y) =
\chi(\cO_{\widetilde Y}) = 1 = \kcat(\widetilde Y^{\vee})$, as each
quantity is a birational invariant of smooth CY$_{3}$'s
(Batyrev--Kontsevich motivic integration).

\subsection*{Cycle 4 --- Kontsevich--Soibelman wall-crossing at the
             character level}

\ClaimStatusTheorem\ Pass to the Poincar\'e character
$\chi_{\Yplus}(q,y_{0},y_{1}) = \sum_{n_{0},n_{1}}
q^{n_{0}+n_{1}}y_{0}^{n_{0}}y_{1}^{n_{1}}\dim\Yplus_{(n_{0},n_{1})}$.
The generating function is the refined DT partition function
$Z^{\mathrm{DT}}_{\widetilde Y}(q,y_{0},y_{1})$ on the NCDT chamber
(Szendr\H{o}i \texttt{0705.3419}, Thm.~1.2).  Kontsevich--Soibelman
(\texttt{0811.2435}, sec.~8.4) identifies $\Wflopp$ at the character
level with the \emph{wall-crossing operator} between the NCDT chamber
and the DT chamber of $\widetilde Y^{\vee}$:
\[
  Z^{\mathrm{NCDT}}_{\widetilde Y}(q,y_{0},y_{1}) \;=\;
  Z^{\mathrm{DT}}_{\widetilde{Y}^{\vee}}(q,y_{1},y_{0})\cdot
  M(-q)^{-2},
\]
where $M(q)=\prod_{n\ge 1}(1-q^{n})^{-n}$ is the MacMahon function and
the factor $M(-q)^{-2}$ is the contribution of the two bifundamental
hypermultiplets crossing the wall (Nagao--Nakajima \texttt{0809.2992},
Thm.~1.1, for refined version).  The exchange $y_{0}
\leftrightarrow y_{1}$ is the image of $s_{1}\in W(A_{1}^{(1)})$ on the
character lattice; the $M(-q)^{-2}$ prefactor is the imaginary-root
parity flip of Cycle 3, evaluated at the universal character.

\subsection*{Cycle 5 --- the Seiberg-duality groupoid}

\ClaimStatusTheorem\ Wave-17 C1, Cycle 4 produces the overlap cocycle
$\phi_{UV}\colon\Yplus(\widetilde Y)|_{U_{+}\cap U_{-}}^{+}
\xrightarrow{\sim}\Yplus(\widetilde Y)|_{U_{+}\cap U_{-}}^{-}$; the
present cycle produces the flop automorphism $\Wflopp$.  The pair
generates a subgroup of $\Aut_{\mathrm{alg}}(\Yplus(\widetilde Y))$
isomorphic to the \emph{Seiberg-duality groupoid} of
$(Q_{\mathrm{KW}},W_{\mathrm{KW}})$: a groupoid whose objects are
presentations of the NCCR $\Jac(Q_{\mathrm{KW}},W_{\mathrm{KW}})$ via
tilting bundles (Van den Bergh \texttt{math/0207170}) and whose
morphisms are mutations at vertices (Aspinwall--Katz
\texttt{hep-th/0412209}).  Explicitly:
\[
  \Wflopp = s_{1}\cdot\pi,\qquad
  \phi_{UV} = s_{0}\cdot s_{1}\cdot\pi,
\]
where $s_{0},s_{1}\in W(A_{1}^{(1)})$ are the simple reflections and
$\pi$ is the Cartan parity flip; their product generates $W(A_{1}^{(1)})
\rtimes\Z/2$, i.e.\ the affine Weyl group extended by the diagram
automorphism.  Mathematically, $\pi^{2}=1,s_{i}^{2}=1$, and
$(s_{0}s_{1})^{n}$ is of infinite order (exactly the affine
$A_{1}^{(1)}$ Weyl group).  The groupoid closes: all presentations of
$\Jac(Q_{\mathrm{KW}},W_{\mathrm{KW}})$ related by iterated mutation are
related, at the level of $\Yplus$, by compositions of $\Wflopp$ and
$\phi_{UV}$.

\emph{Verification}.  Three independent checks.
(i)~Action on $A_{1}^{(1)}$ simple roots: $s_{0}$ sends $\alpha_{0}
\mapsto -\alpha_{0}$, $\alpha_{1}\mapsto\alpha_{0}+\alpha_{1}$; $s_{1}$
sends $\alpha_{1}\mapsto-\alpha_{1},\alpha_{0}\mapsto\alpha_{0}+\alpha_{1}$;
their composition has infinite order and preserves the imaginary root
$\delta=\alpha_{0}+\alpha_{1}$.  (ii)~Action on refined DT generating
series: $\phi_{UV}$ exchanges $(y_{0},y_{1})\mapsto(y_{1},y_{0})$ at
the character level (overlap cocycle, Wave-17 C1 Cycle 4), while
$\Wflopp$ sends $y_{1}\mapsto y_{1}^{-1}q^{?}$ up to the DT
$M(-q)^{-2}$ prefactor (Cycle 4 above); both are well-defined
substitutions preserving the partition function modulo MacMahon.
(iii)~Preservation of $\kch=1$, $\kcat=1$, and the Borcherds weight
$\kBKM=0$ (no modular witness for non-compact CY$_{3}$: Gritsenko
$\Phi_{N}$ is trivial when $\chi$ vanishes at the fibre, and
$\widetilde Y$ is non-compact, so $\kBKM(\widetilde Y) = 0$ by
convention; the flop preserves this trivially).  All three checks
confirm $\Wflopp\in\Aut_{\mathrm{alg}}(\Yplus(\widetilde Y))$.\qed
